\section{Introduction}
\label{sec:intro}

Natural language processing and computer vision benefit from large public corpora such as The Pile
and ImageNet~\cite{gao2021pile,deng2009imagenet}. In integrated-circuit (IC) design, public data that
combine diverse designs with trustworthy physical annotations remain scarce: existing collections
cover particular design sets, implementation stages, or
technologies~\cite{chowdhury2021openabcd,chai2022circuitnet,shi2025forgeeda,shrestha2026edaschema,zhou2026r2g}.
Building reliable datasets is therefore a foundation for training and evaluating larger circuit
models. Circuits are naturally graphs of cells, pins, and nets, and graph representations are widely
used in machine learning for electronic design automation
(ML-for-EDA)~\cite{chowdhury2021openabcd,liang2023circuitops,zhou2026r2g}. The challenge is to produce
those graphs at scale while preserving the circuit identities and physical labels needed for
learning.

\begin{figure*}[t]
  \centering
  \PH{Figure 1 from the existing draft (RTL-to-graph workflow: nine stages from RTL acquisition to ML
  dataset, frozen-recipe repair path, and related-work coverage); import fig/pipeline.pdf}
  \caption{RTL-to-graph workflow and representative prior-work coverage. The nine stages span
  qualified RTL acquisition, physical execution and repair, and audited graph production. The dashed
  path shows reuse of frozen repair recipes after failed signoff checks.}
  \label{fig:overview}
\end{figure*}

\cref{fig:overview} shows why RTL-to-graph dataset construction is difficult: source acquisition and
synthesis feed floorplanning, placement, clock-tree synthesis, routing, signoff, graph construction,
and dataset admission. This chain spans multiple EDA tools, and failures or representation changes at
one stage can invalidate later data. Even before the first design runs, the toolchain itself must be
installed and version-matched. LLM-based agents can automate tool use and physical-flow
decisions~\cite{he2024chateda,ghose2026orfsagent,deng2026rtlgds}; HighTide extends agent-assisted
RTL-to-GDS curation~\cite{goldblatt2026hightide}. Dataset efforts provide synthesis graphs,
physical-design data, and stage-resolved
representations~\cite{chowdhury2021openabcd,chai2022circuitnet,shi2025forgeeda,shrestha2026edaschema,zhou2026r2g}.
As the prior-work coverage in the figure illustrates, these efforts address portions of the chain,
but leave its source-to-graph yield and the semantic validity of the final training examples
insufficiently measured.

We introduce \R, an evidence-gated RTL-to-graph workflow that connects environment provisioning,
qualified source acquisition, physical implementation, and audited graph production. Our
contributions are:
\begin{enumerate}
  \item A workflow, packaged as agent skills and executed by deterministic programs, that provisions
        its own toolchain, builds a graph dataset from public RTL repositories through source
        qualification, physical implementation, and graph construction, records acceptance and loss
        at each handoff, and admits samples only after violation-free physical signoff
        (\cref{sec:system}).
  \item An evaluation of each handoff under fixed protocols with deterministic verdicts and full cost
        accounting: toolchain installation in a clean environment (\cref{sec:exp1}), RTL acquisition
        yield (\cref{sec:exp2}), and signoff repair of 676 failing designs on three process kits,
        comparing the same models with and without \R's validated repair knowledge
        (\cref{sec:exp3}).
  \item Stage-aware graph datasets with independent checks of circuit identity, feature visibility,
        and physical labels, and learning baselines on them (\cref{sec:exp4}).
\end{enumerate}
Together, these contributions make dataset yield and sample correctness explicit objectives.
